\scanpage{108}
\mspage{4}
% archive foliation: 107

en corr. 1-1 avec les éléments de $(t_E^{\otimes2})^*$.

\uline{Corollaire 3} On a
\[
\begin{aligned}
(31)\qquad \mathcal F_\infty&\simeq k[P,P']/P'^2-(P^3+a_1P^2+a_2P^{\unsure{2}}+a_3)\\
&\simeq k[P][P']/(P'^2-(P^3+a_1P^2+a_2\unsure{P^2|P}+a_3))
\end{aligned}
\]
L'isom étant donné par $P\longmapsto\wp$, $P'\longmapsto\wp'$.
L'autom $f\longrightarrow\check f$ correspond à l'autom.
$P'\longmapsto-P'$, $P\longmapsto P$ de $k[P,P']/\ldots$.
La dérivation $\Theta_\delta$ est déduite par passage aux quotients de
\deleted{l'opération} dérivation
\[
(32)\qquad \Theta_\delta P=P',\qquad \Theta P'=\tfrac12(3P^2+2a_1P+a_2)
\]
sur $k[P,P']$, qui applique l'idéal $P'^2-(P^3+a_1P^2+a_2P+a_3)$ en lui-même.

La formule (32) s'obtient en dérivant
$\wp'^2=\wp^3+a_1\wp^2+a_2\wp+a_3$ par $\Theta_\delta$,
\deleted{et en divisant} \inserted{d'où}
\[
2\wp'\Theta_\delta(\wp')=(3\wp^2+2a_1\wp+a_2)\wp'
\]
et en divisant par $\wp'$ la relation obtenue, d'où
\[
(33)\qquad \wp^{\unsure{'|''}}=\tfrac12(3\wp^2+2a_1\wp+a_2).
\]

\deleted{\uline{Corollaire}} \uline{Remarque} En termes des bases explicites plus haut des
$\mathcal F_n^+$ et $\mathcal F_n^-$ (25') donc de
$\mathcal F_n=\mathcal F_n^+\oplus\mathcal F_n^-$, on voit maintenant comment
expliciter la loi de multiplication \inserted{en termes de $a_1,a_2,a_3$,}
ainsi que
